\documentclass[11pt,a4paper,fontset=fandol]{ctexart}

\usepackage[a4paper,top=2.4cm,bottom=2.6cm,left=2.35cm,right=2.35cm,headheight=18pt,headsep=0.55cm,footskip=26pt]{geometry}
\usepackage{amsmath,amssymb}
\usepackage{fancyhdr}
\usepackage{lastpage}
\usepackage{enumitem}
\usepackage{titlesec}
\usepackage{xcolor}
\usepackage{hyperref}
\usepackage{bookmark}
\usepackage[most]{tcolorbox}
\usepackage{setspace}
\usepackage{array}

\hypersetup{
  colorlinks=true,
  linkcolor=blue!50!black,
  urlcolor=blue!50!black
}

\setstretch{1.16}
\setlength{\parindent}{2em}
\setlength{\parskip}{0.32em}
\setlist[enumerate]{itemsep=0.35em, topsep=0.35em, leftmargin=2.3em}
\setlist[itemize]{itemsep=0.25em, topsep=0.25em, leftmargin=2em}

\definecolor{mainblue}{RGB}{36,83,140}
\definecolor{lightblue}{RGB}{241,246,252}
\definecolor{lightgray}{RGB}{246,246,246}
\definecolor{rulegray}{RGB}{110,118,128}
\definecolor{softgreen}{RGB}{236,247,240}
\definecolor{softyellow}{RGB}{254,249,229}

\titleformat{\section}
  {\Large\bfseries\color{mainblue}}
  {\thesection.}
  {0.45em}
  {}
  [\vspace{0.25em}\color{rulegray}\titlerule]
\titlespacing*{\section}{0pt}{1.2em}{0.8em}
\titleformat{\subsection}{\large\bfseries\color{mainblue}}{\thesubsection}{0.5em}{}

\pagestyle{fancy}
\fancyhf{}
\fancyhead[L]{\small 递推计数专题目录页}
\fancyhead[C]{\small 组合专题资料索引}
\fancyhead[R]{\small 刘文博}
\fancyfoot[C]{\small 第 \thepage 页 / 共 \pageref{LastPage} 页}
\renewcommand{\headrulewidth}{0.55pt}
\renewcommand{\footrulewidth}{0.35pt}

\newtcolorbox{goalbox}[1]{
  breakable,
  colback=lightblue,
  colframe=mainblue,
  boxrule=0.7pt,
  arc=1mm,
  title=\textbf{#1},
  fonttitle=\bfseries
}

\newtcolorbox{infobox}[1]{
  breakable,
  colback=softgreen,
  colframe=green!40!black,
  boxrule=0.65pt,
  arc=1mm,
  title=\textbf{#1},
  fonttitle=\bfseries
}

\newtcolorbox{warnbox}[1]{
  breakable,
  colback=softyellow,
  colframe=orange!60!black,
  boxrule=0.65pt,
  arc=1mm,
  title=\textbf{#1},
  fonttitle=\bfseries
}

\begin{document}

\thispagestyle{empty}
\begin{titlepage}
\centering
\vspace*{0.9cm}

\begin{tcolorbox}[
  enhanced,
  width=0.92\textwidth,
  colback=white,
  colframe=mainblue,
  boxrule=1pt,
  arc=0mm,
  left=6mm,
  right=6mm,
  top=8mm,
  bottom=8mm
]
\centering

{\fontsize{24}{30}\selectfont\bfseries 递推计数专题目录页\par}
\vspace{0.65cm}
{\fontsize{17}{22}\selectfont 讲义、题单与周测的使用说明\par}

\vspace{1cm}
\begin{tcolorbox}[colback=lightblue,colframe=mainblue,width=0.84\textwidth,boxrule=1pt]
\centering
{\large 适合对象：已经掌握基础计数，准备系统提升“按最后一步分类”能力的学生}\\[0.35em]
{\normalsize 本目录页帮助把递推计数专题资料串成一个完整训练包}
\end{tcolorbox}

\vspace{1.0cm}
\renewcommand{\arraystretch}{1.42}
\begin{tabular}{>{\bfseries}p{3.5cm}p{8cm}}
专题名称： & 递推计数 \\
资料组成： & 课堂讲义、提高训练题单、周测 A 卷、周测 B 卷 \\
使用方式： & 先学讲义，再做题单，最后用 A/B 周测检查掌握情况 \\
编译方式： & XeLaTeX \\
\end{tabular}

\vspace{0.9cm}
{\large \today\par}
\end{tcolorbox}
\end{titlepage}

\tableofcontents
\newpage

\section{专题包包含哪些资料}

\begin{goalbox}{四份核心资料}
\begin{enumerate}
  \item 课堂讲义：帮助理解递推计数的基本思想，掌握如何按最后一步、最后一段或最后一个位置来分类；
  \item 提高训练题单：用于课后巩固和变式提升，重点训练“会写递推、会设状态、会写初始条件”；
  \item 周测 A 卷：检查基础模型是否稳定，例如上楼梯、受限字符串、简单铺砖；
  \item 周测 B 卷：检查状态设计、综合变式、以及与错位排列等旧专题的联系是否真正掌握。
\end{enumerate}
\end{goalbox}

\begin{infobox}{对应文件名}
\begin{itemize}
  \item 讲义：\texttt{recursive\_counting\_handout.tex} / \texttt{recursive\_counting\_handout.pdf}
  \item 提高训练题单：\texttt{recursive\_counting\_advanced\_problem\_set.tex} / \texttt{recursive\_counting\_advanced\_problem\_set.pdf}
  \item 周测 A 卷：\texttt{recursive\_counting\_weekly\_test\_A.tex} / \texttt{recursive\_counting\_weekly\_test\_A.pdf}
  \item 周测 B 卷：\texttt{recursive\_counting\_weekly\_test\_B.tex} / \texttt{recursive\_counting\_weekly\_test\_B.pdf}
\end{itemize}
\end{infobox}

\section{建议学习顺序}

\begin{goalbox}{推荐顺序}
\begin{enumerate}
  \item 先完整学习讲义，重点吃透“为什么按最后一步分类”“什么时候要增加状态”；
  \item 再完成提高训练题单，把基本模型和变式模型做熟；
  \item 接着做周测 A 卷，检查基础递推是否稳定；
  \item 最后做周测 B 卷，检查综合迁移、状态设计和与旧专题的联系。
\end{enumerate}
\end{goalbox}

\begin{warnbox}{不建议的做法}
\begin{itemize}
  \item 只背递推式，不写清楚它是怎么分类得到的；
  \item 只写递推，不写初始条件；
  \item 分类时既不完整，也不检查是否有重叠；
  \item 做完题只看最后结果，不回头检查“状态设得是否合适”。
\end{itemize}
\end{warnbox}

\section{每份资料主要解决什么问题}

\subsection{课堂讲义}
讲义适合第一次系统学习这个专题，主要解决以下问题：
\begin{itemize}
  \item 什么叫递推计数；
  \item 哪些题适合用递推而不是直接计数；
  \item 如何从“最后一步”“最后一段”“最后一块”切入；
  \item 什么情况下需要增加状态；
  \item 递推计数与错位排列、受限字符串、铺砖模型之间的联系。
\end{itemize}

\subsection{提高训练题单}
题单适合第二阶段训练，主要用来把方法做熟。它更强调：
\begin{itemize}
  \item 能否独立建立递推关系；
  \item 能否正确写出初始条件；
  \item 能否识别“需要增加状态”的题；
  \item 能否处理铺砖、字符串、楼梯等不同模型之间的迁移。
\end{itemize}

\subsection{周测 A 卷}
A 卷偏基础，适合检查：
\begin{itemize}
  \item 斐波那契型递推是否熟练；
  \item 基础上楼梯模型是否会做；
  \item 不含相邻限制的字符串与铺砖模型是否会做；
  \item 错位排列递推是否能顺利接上旧专题。
\end{itemize}

\subsection{周测 B 卷}
B 卷偏综合，适合检查：
\begin{itemize}
  \item 是否会用两个状态或多个状态；
  \item 是否会从递推中看出不同模型的联系；
  \item 是否会处理更灵活的字符串与铺砖变式；
  \item 是否会解释“为什么这个递推成立”。
\end{itemize}

\section{一个推荐的 4 次训练安排}

\begin{goalbox}{4 次完成一个小专题}
\begin{enumerate}
  \item 第 1 次：学习讲义前半部分，重点理解递推思想与基本模型；
  \item 第 2 次：学习讲义后半部分，完成题单基础题和中档题；
  \item 第 3 次：完成题单综合题，并把错题按“分类漏了/重了/初值错了/状态不够”分类订正；
  \item 第 4 次：完成周测 A 卷或 B 卷，作为阶段性检查。
\end{enumerate}
\end{goalbox}

\begin{infobox}{如果孩子基础较好，可以这样压缩}
\begin{itemize}
  \item 第 1 天：讲义 + 题单基础题；
  \item 第 2 天：题单变式题 + 错题订正；
  \item 第 3 天：周测 A 卷；
  \item 第 4 天：周测 B 卷。
\end{itemize}
\end{infobox}

\section{专题学习后的自检清单}

\begin{goalbox}{如果下面 8 条都能做到，说明这个专题基本过关}
\begin{enumerate}
  \item 我能说清楚递推计数为什么是“把大问题拆成小问题”；
  \item 我会按最后一步或最后一段分类；
  \item 我不会漏写初始条件；
  \item 我知道什么时候需要增加状态；
  \item 我能处理不含相邻限制的 0,1 串；
  \item 我能处理基本铺砖模型；
  \item 我知道递推计数和错位排列递推之间的联系；
  \item 我会在算完以后检查递推式、初值和结果是否合理。
\end{enumerate}
\end{goalbox}

\section{给家长的使用建议}

\begin{warnbox}{更有效的带练方式}
\begin{itemize}
  \item 不要一开始就提示公式，先让孩子说“最后一步可能是什么”；
  \item 如果卡住，可以只提示“分最后一块试试”或“是不是要多设一个状态”；
  \item 订正时重点看递推关系是怎么来的，而不只是最后数值是否正确；
  \item 错题建议按错误类型整理：递推式错、初值错、分类不完整、状态不够。
\end{itemize}
\end{warnbox}

\begin{infobox}{和后续专题的衔接}
递推计数学完以后，可以自然过渡到更深入的\textbf{状态压缩递推}、\textbf{生成函数初步}，也可以回头和\textbf{容斥原理}、\textbf{错位排列}互相对照，体会同一问题可以有不同解法。
\end{infobox}

\end{document}
